\documentclass[11pt]{article}

\usepackage[margin=1in]{geometry}
\usepackage{amsmath}
\usepackage{booktabs}
\usepackage{graphicx}
\usepackage{natbib}
\usepackage[hidelinks]{hyperref}

\title{Educational Upgrading and Labor Income Growth in Latin America}
\author{Oliver Pardo}
\date{\today}

\begin{document}

\maketitle

\begin{abstract}
This paper decomposes changes in average monthly labor income into changes in
the educational composition of employment and changes in income within
education groups. For each of 14 Latin American economies in the World Bank's
LAC Equity Lab, we compare the oldest and newest available fourth-quarter
observations. The endpoints usually correspond to 2016 and 2023, although
coverage differs by economy. The education component is positive in 12 of the
14 economies, while within-group income changes are negative in 10. Synthetic
mean income increases in 8 economies and falls in 6. These are country-level
accounting decompositions, not causal estimates of the return to education.
Four comparisons cross documented survey-series changes and should be
interpreted with additional caution.
\end{abstract}

\section{Introduction}

Average labor income can change for two distinct reasons. First, the
educational composition of employment can change. Second, the mean income
received by workers with a given educational attainment can rise or fall.
The first channel measures educational upgrading; the second measures changes
within education groups. Distinguishing the two is useful because an increase
in average labor income need not imply that workers at each educational level
are receiving higher income.

This paper measures both contributions for 14 Latin American economies
covered by the World Bank's LAC Equity Lab. For each economy, the comparison
uses the oldest and newest available fourth-quarter observation. Restricting
the analysis to the fourth quarter holds the observed season fixed and
prevents differences in quarterly coverage from entering the decomposition.
The approach also gives each economy its longest available comparison instead
of imposing a common window that would exclude countries.

The decomposition follows the symmetric method associated with
\citet{kitagawa1955} and the broader decomposition literature reviewed by
\citet{fortinlemieuxfirpo2011}. The exercise is related to work documenting
educational upgrading and changing skill premia in Latin America
\citep{gasparinietal2011,azevedoetal2013,messinadasilva2019}. Its contribution
is narrower: it applies one exact accounting identity to the longest
fourth-quarter comparison available for each economy and reports both the
underlying descriptive statistics and the country-level results.

\section{Descriptive statistics}

\subsection{Data and educational groups}

The LAC Equity Lab reports mean monthly labor income and the number of workers
by educational attainment \citep{lablac}. Monthly labor income is expressed in
2017 purchasing-power-parity dollars and covers employed people aged 15 or
older with coherent reported income. The worker counts refer to the employed
population, not to the labor force, which would also include unemployed
people.

LABLAC harmonizes educational attainment into three groups. \emph{Low}
education includes people who never attended school, completed primary
education, or did not complete secondary education. \emph{Middle} education
includes people who completed secondary education or began but did not
complete higher education. \emph{High} education includes people who
completed higher education. The three groups are used consistently in the
income and worker-count tables.

The downloaded data contain Argentina, Bolivia, Brazil, Chile, Colombia,
Costa Rica, the Dominican Republic, Ecuador, El Salvador, Guatemala, Mexico,
Paraguay, Peru (Lima and Callao), and Uruguay. Peru is a subnational
observation and should not be interpreted as a national estimate. The initial
year is 2016 for every economy except Paraguay, whose first usable fourth
quarter is 2017. The final year is 2023 for most economies, 2024 for Costa
Rica and Ecuador, and 2022 for Guatemala.

\subsection{Income levels}

Table~\ref{tab:latam_country_income_descriptives} reports two income measures.
The published total is LABLAC's mean for all workers with coherent reported
income. The synthetic mean weights the three education-group income means
with their shares of published workers. The decomposition uses the synthetic
mean because it is exactly consistent with the observed education shares.

\begin{table}[htbp]
\centering
\caption{Endpoint coverage and monthly labor income}
\label{tab:latam_country_income_descriptives}
\small
\resizebox{\textwidth}{!}{%
\begin{tabular}{lrrrrrr}
\toprule
Economy & Initial & Final & Published initial & Published final & Synthetic initial & Synthetic final \\
\midrule
Argentina & 2016 & 2023 & 1,267 & 918 & 1,267 & 918 \\
Bolivia & 2016 & 2023 & 997 & 958 & 980 & 943 \\
Brazil & 2016 & 2023 & 897 & 945 & 892 & 942 \\
Chile & 2016 & 2023 & 1,174 & 1,291 & 1,160 & 1,278 \\
Colombia & 2016 & 2023 & 755 & 908 & 749 & 903 \\
Costa Rica & 2016 & 2024 & 1,328 & 1,315 & 1,298 & 1,304 \\
Dominican Republic & 2016 & 2023 & 800 & 837 & 799 & 835 \\
Ecuador & 2016 & 2024 & 843 & 708 & 838 & 694 \\
El Salvador & 2016 & 2023 & 564 & 690 & 559 & 687 \\
Guatemala & 2016 & 2022 & 614 & 634 & 609 & 629 \\
Mexico & 2016 & 2023 & 614 & 673 & 640 & 698 \\
Paraguay & 2017 & 2023 & 1,007 & 1,002 & 999 & 994 \\
Peru (Lima and Callao) & 2016 & 2023 & 927 & 829 & 925 & 825 \\
Uruguay & 2016 & 2023 & 1,130 & 1,129 & 1,129 & 1,128 \\
\bottomrule
\end{tabular}
}
\begin{minipage}{0.98\textwidth}
\footnotesize
\textit{Note:} All observations are fourth-quarter values in 2017 PPP dollars. Published income is LABLAC's reported total. Synthetic income weights the three education-group means with the published number of workers in each group.
\end{minipage}
\end{table}


The two measures are close but are not identical. Across the 28 selected
endpoints, the largest absolute difference between synthetic and published
income is 4.2 percent. The difference arises because the education-group
worker counts include all workers in each group, whereas the income means use
workers with coherent reported income. The synthetic measure therefore
defines the estimand; the published total is a diagnostic.

\subsection{Educational composition and group income}

Table~\ref{tab:latam_country_education_shares} reports the observed composition
of employment. The share with low education falls in 11 of the 14 economies.
The largest reductions occur in Chile, Colombia, Bolivia, Brazil, Costa Rica,
and Paraguay. Ecuador and Guatemala move in the opposite direction: the low
education share rises by 4.0 and 7.0 percentage points, respectively.

\begin{table}[htbp]
\centering
\caption{Composition of employment by educational attainment}
\label{tab:latam_country_education_shares}
\small
\resizebox{\textwidth}{!}{%
\begin{tabular}{lrrrrrr}
\toprule
Economy & Low initial & Low final & Middle initial & Middle final & High initial & High final \\
\midrule
Argentina & 36.0 & 29.9 & 41.0 & 44.6 & 22.9 & 25.5 \\
Bolivia & 57.4 & 47.4 & 36.5 & 44.9 & 6.1 & 7.7 \\
Brazil & 42.6 & 33.0 & 40.5 & 45.3 & 16.9 & 21.6 \\
Chile & 31.9 & 19.8 & 53.9 & 55.3 & 14.1 & 24.8 \\
Colombia & 46.1 & 34.9 & 40.2 & 35.6 & 13.7 & 29.6 \\
Costa Rica & 54.5 & 44.9 & 32.3 & 38.6 & 13.1 & 16.5 \\
Dominican Republic & 55.3 & 48.4 & 31.4 & 36.6 & 13.4 & 15.0 \\
Ecuador & 44.4 & 48.4 & 38.1 & 36.6 & 17.5 & 15.0 \\
El Salvador & 75.2 & 72.6 & 20.6 & 21.0 & 4.2 & 6.4 \\
Guatemala & 72.1 & 79.1 & 23.1 & 15.0 & 4.7 & 5.9 \\
Mexico & 61.7 & 53.3 & 21.6 & 25.5 & 16.7 & 21.2 \\
Paraguay & 56.1 & 46.7 & 31.3 & 36.0 & 12.6 & 17.3 \\
Peru (Lima and Callao) & 17.5 & 16.0 & 51.4 & 49.4 & 31.2 & 34.6 \\
Uruguay & 64.1 & 60.3 & 22.3 & 22.9 & 13.6 & 16.8 \\
\bottomrule
\end{tabular}
}
\begin{minipage}{0.98\textwidth}
\footnotesize
\textit{Note:} Percent of employed workers in each education group. Initial and final refer to the years reported in Table~\ref{tab:latam_country_income_descriptives}.
\end{minipage}
\end{table}


The high education share rises in 13 economies. Colombia records the largest
increase, 15.8 percentage points, followed by Chile with 10.7 points. Ecuador
is the only economy in which the high education share falls. Colombia and
Ecuador also cross documented series changes, so those movements combine
changes in the employed population with possible effects from survey
comparability.

Table~\ref{tab:latam_country_education_incomes} reports mean monthly labor
income for each education group. High-education workers receive more than
middle- and low-education workers in every selected endpoint. The table also
shows that educational upgrading can coexist with falling income within the
education groups. That distinction motivates the decomposition.

\begin{table}[htbp]
\centering
\caption{Monthly labor income by educational attainment}
\label{tab:latam_country_education_incomes}
\small
\resizebox{\textwidth}{!}{%
\begin{tabular}{lrrrrrr}
\toprule
Economy & Low initial & Low final & Middle initial & Middle final & High initial & High final \\
\midrule
Argentina & 863 & 622 & 1,214 & 834 & 1,995 & 1,410 \\
Bolivia & 801 & 760 & 1,159 & 1,031 & 1,588 & 1,549 \\
Brazil & 523 & 535 & 767 & 759 & 2,125 & 1,944 \\
Chile & 632 & 674 & 1,080 & 1,019 & 2,660 & 2,337 \\
Colombia & 459 & 440 & 688 & 659 & 1,897 & 1,743 \\
Costa Rica & 855 & 838 & 1,374 & 1,259 & 2,951 & 2,681 \\
Dominican Republic & 640 & 648 & 726 & 797 & 1,628 & 1,533 \\
Ecuador & 577 & 501 & 790 & 713 & 1,604 & 1,269 \\
El Salvador & 465 & 587 & 773 & 849 & 1,208 & 1,303 \\
Guatemala & 454 & 535 & 907 & 763 & 1,526 & 1,556 \\
Mexico & 492 & 544 & 652 & 668 & 1,174 & 1,122 \\
Paraguay & 751 & 787 & 1,062 & 995 & 1,950 & 1,554 \\
Peru (Lima and Callao) & 548 & 485 & 699 & 644 & 1,507 & 1,240 \\
Uruguay & 861 & 786 & 1,237 & 1,204 & 2,216 & 2,254 \\
\bottomrule
\end{tabular}
}
\begin{minipage}{0.98\textwidth}
\footnotesize
\textit{Note:} Mean monthly labor income in 2017 PPP dollars. All observations correspond to the fourth quarter.
\end{minipage}
\end{table}


\section{Decomposition}

\subsection{Endpoint construction and weights}

For each economy, the calculation selects the oldest and newest available
fourth-quarter observation with positive income and worker counts for all
three education groups. If more than one Tableau source is present in an
endpoint, the calculation selects the complete source that most closely
reproduces the values displayed in the Equity Lab workbook. It then uses the
worker counts reported for that same source.

The education weights are observed rather than inferred from the published
total income. With three education groups, the published total supplies the
identity
\[
  R_t=s_{L,t}r_{L,t}+s_{M,t}r_{M,t}+s_{H,t}r_{H,t},
\]
while the shares also satisfy
\[
  s_{L,t}+s_{M,t}+s_{H,t}=1.
\]
These are only two independent equations for three unknown shares. Therefore,
the three weights cannot be recovered uniquely from total income and the
three group means. We instead calculate each share directly from LABLAC's
reported worker counts:
\[
  s_{g,t}=\frac{N_{g,t}}{\sum_j N_{j,t}}.
\]

\subsection{Exact decomposition}

Let \(R_t\) denote average remuneration in period \(t\). The population and
remuneration concept are defined separately in each paper. Educational
attainment partitions that population into mutually exclusive groups indexed by
\(g\):
\begin{equation}
  R_t = \sum_g s_{g,t} r_{g,t},
  \label{eq:remuneration_identity}
\end{equation}
where \(s_{g,t}\) is group \(g\)'s share of the relevant population and
\(r_{g,t}\) is its average remuneration.

For two periods, the change in average remuneration can be written as
\begin{equation}
  R_1-R_0 =
  \sum_g \bar{s}_g\left(r_{g,1}-r_{g,0}\right)
  +
  \sum_g \bar{r}_g\left(s_{g,1}-s_{g,0}\right),
  \label{eq:symmetric_decomposition}
\end{equation}
where
\(\bar{s}_g=(s_{g,1}+s_{g,0})/2\) and
\(\bar{r}_g=(r_{g,1}+r_{g,0})/2\).
The first term measures changes in remuneration within educational groups. The
second measures changes in the educational composition of the relevant
population. This symmetric decomposition is exactly additive and does not
depend on choosing either period as the reference.

Equation~\eqref{eq:symmetric_decomposition} is an accounting identity. The
within-group component is not a measure of productivity, and the composition
component does not identify the causal effect of education. Changes in either
component may reflect selection, labor supply, labor demand, institutions,
hours, occupations, sectors, or survey coverage.


The within-group term in
Equation~\eqref{eq:symmetric_decomposition} weights the change in each group's
income by its average employment share. The education term weights the change
in each group's employment share by its average income. Both components are
reported as percentages of initial synthetic income. They add exactly to the
percentage change in synthetic income before rounding.

\subsection{Interpretation}

The decomposition is an accounting identity. A positive education component
means that employment shifted toward education groups with higher average
income. It does not identify the causal return to schooling. A positive
within-group component means that income rose, on average, within the
education groups after holding their employment shares fixed. It need not
measure productivity because income within a group can also change with
hours, occupations, sectors, formality, labor demand, selection into
employment, and survey coverage.

\section{Results}

Table~\ref{tab:latam_country_decomposition} reports one decomposition for each
economy. Synthetic mean income rises in 8 economies and falls in 6. El
Salvador records the largest increase, 22.9 percent, followed by Colombia with
20.7 percent. Argentina records the largest decline, 27.6 percent, followed by
Ecuador with 17.2 percent and Peru (Lima and Callao) with 10.8 percent.

\begin{table}[htbp]
\centering
\caption{Country-level decomposition of synthetic labor income}
\label{tab:latam_country_decomposition}
\small
\resizebox{\textwidth}{!}{%
\begin{tabular}{lrrrrrr}
\toprule
Economy & Initial & Final & Total change & Education & Within groups & Same source \\
\midrule
Argentina & 2016 & 2023 & -27.6 & 2.8 & -30.3 & Yes \\
Bolivia & 2016 & 2023 & -3.8 & 4.0 & -7.8 & Yes \\
Brazil & 2016 & 2023 & 5.6 & 9.3 & -3.7 & Yes \\
Chile & 2016 & 2023 & 10.1 & 17.5 & -7.4 & Yes \\
Colombia & 2016 & 2023 & 20.7 & 27.6 & -6.9 & No \\
Costa Rica & 2016 & 2024 & 0.5 & 7.3 & -6.9 & Yes \\
Dominican Republic & 2016 & 2023 & 4.5 & 2.7 & 1.9 & Yes \\
Ecuador & 2016 & 2024 & -17.2 & -3.1 & -14.2 & No \\
El Salvador & 2016 & 2023 & 22.9 & 3.1 & 19.8 & Yes \\
Guatemala & 2016 & 2022 & 3.3 & -2.5 & 5.8 & Yes \\
Mexico & 2016 & 2023 & 9.1 & 5.3 & 3.8 & Yes \\
Paraguay & 2017 & 2023 & -0.5 & 5.8 & -6.4 & No \\
Peru (Lima and Callao) & 2016 & 2023 & -10.8 & 2.9 & -13.6 & Yes \\
Uruguay & 2016 & 2023 & -0.1 & 4.2 & -4.3 & No \\
\bottomrule
\end{tabular}
}
\begin{minipage}{0.98\textwidth}
\footnotesize
\textit{Note:} Total change, education, and within-group components are percentages of initial synthetic monthly income. The underlying income measure is 2017 PPP US dollars per employed person per month. The two components add to total change before rounding. ``Same source'' indicates that the Tableau series and survey identifiers are unchanged across endpoints.
\end{minipage}
\end{table}


The education component is positive in 12 of the 14 economies. It is largest
in Colombia, where it contributes an amount equivalent to 27.6 percent of
initial synthetic income, and in Chile, where it contributes 17.5 percent.
Brazil follows with a 9.3 percent contribution. In Argentina, Bolivia, Brazil,
Chile, Costa Rica, Paraguay, Peru, and Uruguay, educational upgrading raises
synthetic income but is partly or fully offset by declining income within the
education groups.

The within-group component is negative in 10 economies. Argentina has the
largest decline, equivalent to 30.3 percent of initial synthetic income.
Ecuador and Peru also have large negative within-group components, at 14.2
and 13.6 percent. In contrast, within-group income raises synthetic income in
El Salvador, Guatemala, Mexico, and the Dominican Republic.

Two economies do not exhibit the general pattern of educational upgrading.
In Ecuador, the education component subtracts 3.1 percent of initial income
because employment shifts toward the low education group. In Guatemala, it
subtracts 2.5 percent for the same reason. Guatemala uses the same survey
series at both endpoints. Ecuador does not, so its result cannot be separated
from the documented change in series using these aggregates alone.

Four comparisons cross a Tableau series change: Colombia, Ecuador, Paraguay,
and Uruguay. Their decompositions remain exact descriptions of the published
endpoint aggregates, but part of the measured change may reflect differences
in survey design or processing. Colombia's large positive education
component and Ecuador's negative component therefore require particular
caution. The remaining ten comparisons retain the same survey and series
identifiers at both endpoints.

\section{Conclusion}

The country-level evidence shows a broad increase in the educational
attainment of employed workers. The education component raises synthetic mean
labor income in 12 of 14 economies. However, income within education groups
falls in 10 economies and offsets part of the contribution from educational
upgrading. As a result, synthetic mean income rises in only 8 economies.

The distinction matters for interpreting changes in average labor income.
Educational upgrading can raise the average even when workers within each
education group receive less. Conversely, El Salvador, Mexico, Guatemala,
and the Dominican Republic combine positive within-group income changes with
different movements in educational composition.

These results do not establish why income changed within an education group
or what caused employment to shift across groups. They also do not support a
single aggregate estimate for Latin America because countries have different
endpoint years, Peru covers only Lima and Callao, and four comparisons cross
survey-series changes. The results instead provide a transparent
country-by-country accounting of how educational composition and
within-group income jointly determine the observed change in average labor
income.

\bibliographystyle{apalike}
\bibliography{references}

\end{document}
